% =====================================================================
%  CARNET D'ENTRAINEMENT — FICHE 6 : Liaisons et structures de Lewis
%  THÈME 4 — VSEPR et géométrie moléculaire
%  Fiche TUTO (à lire avant les exercices)
%  Compilation : pdflatex
% =====================================================================
\documentclass[11pt,a4paper]{article}
\usepackage[utf8]{inputenc}
\usepackage[T1]{fontenc}
\usepackage{lmodern}
\usepackage[french]{babel}
\usepackage{amsmath,amssymb,mathtools}
\usepackage{icomma}
\usepackage[version=4]{mhchem}
\usepackage{siunitx}
\sisetup{output-decimal-marker={,}}
\usepackage{xcolor}
\usepackage{colortbl}
\usepackage{tikz}
\usepackage[most]{tcolorbox}
\usepackage{enumitem}
\usepackage{array}
\usepackage{booktabs}
\usepackage{multicol}
\usepackage{titlesec}
\usepackage{fancyhdr}
\usepackage[hidelinks]{hyperref}
\usetikzlibrary{arrows.meta,positioning,calc,shapes.geometric}
\usepackage[a4paper,margin=1.8cm,headheight=26pt,headsep=8pt]{geometry}
\titlespacing*{\section}{0pt}{5pt}{2pt}
\titleformat{\section}{\large\bfseries\color{primarydark}}{}{0pt}{}

% -------- palette --------
\definecolor{primary}{RGB}{146,95,160}
\definecolor{primarydark}{RGB}{92,52,104}
\definecolor{defcol}{RGB}{0,114,178}
\definecolor{thmcol}{RGB}{0,140,120}
\definecolor{excol}{RGB}{0,135,90}
\definecolor{methcol}{RGB}{90,90,100}
\definecolor{attncol}{RGB}{200,40,55}
\definecolor{astucecol}{RGB}{198,93,9}
% -------- atomes --------
\definecolor{atomC}{RGB}{80,80,88}
\definecolor{atomH}{RGB}{235,235,235}
\definecolor{atomO}{RGB}{220,55,45}
\definecolor{atomN}{RGB}{48,80,200}
\definecolor{atomF}{RGB}{150,210,120}
\definecolor{atomCl}{RGB}{70,170,80}
\definecolor{atomS}{RGB}{225,190,40}
\definecolor{atomP}{RGB}{225,140,40}
\definecolor{atomXe}{RGB}{60,180,200}
\definecolor{atomBr}{RGB}{150,70,40}
\definecolor{atomB}{RGB}{230,160,160}
\definecolor{atomSn}{RGB}{120,130,140}
\definecolor{lonepair}{RGB}{120,94,200}
\definecolor{bondcol}{RGB}{60,60,70}

% -------- macros molécules --------
\newcommand{\ball}[5]{\shade[ball color=#2] (#1) circle (#3);\node[font=\footnotesize\bfseries,text=#4] at (#1){#5};}
\newcommand{\pbond}[2]{\draw[bondcol,line width=1.6pt] (#1)--(#2);}
\newcommand{\wbond}[2]{\draw[bondcol,fill=bondcol,draw=none] let \p1=($(#2)-(#1)$), \n1={atan2(\y1,\x1)} in (#1) -- ($(#2)+(\n1+90:0.12)$) -- ($(#2)+(\n1-90:0.12)$) -- cycle;}
\newcommand{\hbond}[2]{\draw[bondcol,line width=1pt,dash pattern=on 1.4pt off 1.8pt] (#1)--(#2);}
\newcommand{\lobe}[3]{\begin{scope}[shift={(#1)},rotate=#2]\shade[ball color=lonepair!55,opacity=0.9] (#3,0) ellipse (0.34 and 0.20);\end{scope}}
\newcommand{\AX}[2]{\ensuremath{\mathrm{AX_{#1}E_{#2}}}}

% -------- boîtes --------
\tcbset{boxbase/.style={enhanced,breakable,boxrule=0.9pt,arc=2.5pt,
   left=3mm,right=3mm,top=2.2mm,bottom=2.2mm,
   fonttitle=\bfseries,coltitle=white,toptitle=0.6mm,bottomtitle=0.6mm}}
\newtcolorbox{defbox}[1][]{boxbase,colback=defcol!6,colframe=defcol,
   title={\textsc{À retenir}\ifx\relax#1\relax\relax\else\textmd{ : \itshape #1}\fi}}
\newtcolorbox{methbox}[1][]{boxbase,colback=methcol!8,colframe=methcol,sharp corners,
   title={\textsc{Méthode}\ifx\relax#1\relax\relax\else\textmd{ : \itshape #1}\fi}}
\newtcolorbox{exbox}[1][]{boxbase,colback=excol!6,colframe=excol,
   title={\textsc{Exemple}\ifx\relax#1\relax\relax\else\textmd{ : \itshape #1}\fi}}
\newtcolorbox{attnbox}[1][]{boxbase,colback=attncol!6,colframe=attncol,
   title={\textsc{Piège à éviter}\ifx\relax#1\relax\relax\else\textmd{ : \itshape #1}\fi}}
\newtcolorbox{astucebox}[1][]{boxbase,colback=astucecol!8,colframe=astucecol,
   title={\textsc{Astuce}\ifx\relax#1\relax\relax\else\textmd{ : \itshape #1}\fi}}

% -------- en-tête --------
\pagestyle{fancy}\fancyhf{}
\fancyhead[L]{\small\textcolor{primarydark}{\textbf{Fiche 6 · Thème 4} — VSEPR et géométrie moléculaire}}
\fancyhead[R]{\small\textcolor{primarydark}{\textsc{Tuto}}}
\fancyfoot[C]{\small\thepage}
\renewcommand{\headrulewidth}{0.8pt}
\renewcommand{\headrule}{\hbox to\headwidth{\color{primary}\leaders\hrule height \headrulewidth\hfill}}

\setlength{\parindent}{0pt}
\setlength{\parskip}{2pt}

\begin{document}

\begin{tcolorbox}[boxbase,colback=primary!12,colframe=primary,
   title={\Large Thème 4 — VSEPR et géométrie moléculaire}]
\textbf{Objectif du thème.} À partir d'une structure de Lewis, \emph{prédire la forme dans
l'espace} d'une molécule. Le modèle \textbf{VSEPR} (Gillespie) repose sur une idée unique : autour
de l'atome central, les doublets d'électrons — liants \emph{et} libres — se repoussent et
s'écartent le plus possible. On code cette information par le symbole \AX{n}{m}, puis on lit la
géométrie dans un tableau.
\end{tcolorbox}

\section*{1. Le code \texorpdfstring{\AX{n}{m}}{AXnEm} : compter les sites de répulsion}

Autour de l'atome central \textbf{A} :
\begin{itemize}[leftmargin=1.6em,topsep=1pt,itemsep=1pt]
  \item \textbf{n} = nombre d'atomes \emph{liés} à A (les \textbf{X}) ;
  \item \textbf{m} = nombre de \textbf{doublets non liants} portés par A (les \textbf{E}).
\end{itemize}
On écrit alors \AX{n}{m}. Le total $n+m$ est le \textbf{nombre de sites de répulsion} : il fixe la
\textbf{géométrie de répulsion} (la façon dont les sites s'écartent). En « masquant » ensuite les
doublets libres, on ne garde que les atomes : c'est la \textbf{géométrie moléculaire}, la forme
réellement observée.

\begin{defbox}[la règle d'or du comptage]
Une \textbf{liaison multiple} (double ou triple) compte pour \textbf{un seul} site de répulsion :
c'est la \emph{direction} de la liaison qui compte, pas le nombre de doublets qu'elle contient.
Ainsi \ce{CO2} (\ce{O=C=O}), avec ses deux doubles liaisons, est du type \AX{2}{0}.
\end{defbox}

\begin{methbox}[trouver la géométrie en 3 gestes]
\begin{enumerate}[leftmargin=1.8em,topsep=1pt,itemsep=1pt]
  \item Écrire la structure de Lewis et repérer l'atome central A.
  \item Compter $n$ (atomes liés) et $m$ (doublets libres sur A) $\Rightarrow$ écrire \AX{n}{m}.
  \item $n+m \to$ \textbf{géométrie de répulsion} ; en masquant les $m$ doublets $\to$
        \textbf{géométrie moléculaire} (lire le tableau ci-dessous).
\end{enumerate}
\end{methbox}

\section*{2. L'exemple canonique : \texorpdfstring{\ce{CH4}}{CH4}, \texorpdfstring{\ce{NH3}}{NH3}, \texorpdfstring{\ce{H2O}}{H2O}}

Ces trois molécules ont \emph{toutes} quatre sites de répulsion (géométrie de répulsion
\textbf{tétraédrique}), mais un nombre croissant de doublets libres : leur géométrie moléculaire
change, et l'angle se \emph{referme} progressivement.

\begin{figure}[h]
\centering
\begin{tikzpicture}[scale=0.92]
  % ===== AX4 : CH4 =====
  \begin{scope}
    \coordinate (C) at (0,0);
    \coordinate (Ha) at (0,1.5);
    \coordinate (Hb) at (-1.35,-0.75);
    \coordinate (Hc) at (1.25,-0.55);
    \coordinate (Hd) at (0.35,-1.35);
    \hbond{C}{Hd}\pbond{C}{Ha}\pbond{C}{Hb}\wbond{C}{Hc}
    \ball{Ha}{atomH}{0.32}{black}{H}
    \ball{Hb}{atomH}{0.32}{black}{H}
    \ball{Hc}{atomH}{0.32}{black}{H}
    \ball{Hd}{atomH}{0.32}{black}{H}
    \ball{C}{atomC}{0.45}{white}{C}
    \node[font=\footnotesize,align=center] at (0,-2.15)
      {\ce{CH4} : \AX{4}{0}\\\textbf{tétraédrique}\\$109{,}5^\circ$};
  \end{scope}
  % ===== AX3E : NH3 =====
  \begin{scope}[xshift=5.2cm]
    \coordinate (N) at (0,0);
    \coordinate (Ha) at (-1.3,-0.85);
    \coordinate (Hb) at (1.2,-0.65);
    \coordinate (Hc) at (0.15,-1.4);
    \pbond{N}{Ha}\wbond{N}{Hb}\hbond{N}{Hc}
    \lobe{N}{90}{0.95}
    \node[lonepair,font=\scriptsize] at (0,1.25) {doublet libre};
    \ball{Ha}{atomH}{0.32}{black}{H}
    \ball{Hb}{atomH}{0.32}{black}{H}
    \ball{Hc}{atomH}{0.32}{black}{H}
    \ball{N}{atomN}{0.45}{white}{N}
    \node[font=\footnotesize,align=center] at (0,-2.15)
      {\ce{NH3} : \AX{3}{1}\\\textbf{pyramide trigonale}\\$\approx 107^\circ$};
  \end{scope}
  % ===== AX2E2 : H2O =====
  \begin{scope}[xshift=10.4cm]
    \coordinate (O) at (0,0);
    \coordinate (Ha) at (-1.3,-0.85);
    \coordinate (Hb) at (1.3,-0.85);
    \pbond{O}{Ha}\pbond{O}{Hb}
    \lobe{O}{68}{0.95}\lobe{O}{112}{0.95}
    \node[lonepair,font=\scriptsize] at (0,1.3) {2 doublets};
    \ball{Ha}{atomH}{0.32}{black}{H}
    \ball{Hb}{atomH}{0.32}{black}{H}
    \ball{O}{atomO}{0.45}{white}{O}
    \node[font=\footnotesize,align=center] at (0,-2.15)
      {\ce{H2O} : \AX{2}{2}\\\textbf{coudée}\\$\approx 104{,}5^\circ$};
  \end{scope}
\end{tikzpicture}
\end{figure}

\section*{3. La hiérarchie des répulsions}

Tous les doublets ne se repoussent pas également. L'ordre décroissant des répulsions est :
$(\text{non liant}\text{--}\text{non liant}) > (\text{non liant}\text{--}\text{liant}) >
(\text{liant}\text{--}\text{liant})$.

\begin{attnbox}[un doublet libre comprime les angles]
Un doublet non liant, plus « volumineux », repousse plus fort que les liaisons : il \textbf{ferme}
les angles. D'où la suite \ce{CH4} $109{,}5^\circ \to$ \ce{NH3} $107^\circ \to$ \ce{H2O}
$104{,}5^\circ$. En \textbf{bipyramide trigonale}, les doublets libres se logent toujours en
position \textbf{équatoriale}, où ils gênent le moins.
\end{attnbox}

\section*{4. Tableau de synthèse VSEPR}

\begin{table}[h]
\centering
\renewcommand{\arraystretch}{1.2}
\setlength{\tabcolsep}{5pt}
\footnotesize
\begin{tabular}{>{\centering\arraybackslash}m{1.2cm} c c >{\raggedright\arraybackslash}m{3.5cm} c c}
\toprule
\rowcolor{primary!22}
\textbf{Sites} & \textbf{Type} & \textbf{Géom. répulsion} & \textbf{Géométrie moléculaire} & \textbf{Angle} & \textbf{Ex.}\\
\midrule
\rowcolor{primary!5} 2 & \AX{2}{0} & linéaire & linéaire & $180^\circ$ & \ce{CO2}\\
3 & \AX{3}{0} & triangulaire & triangulaire plane & $120^\circ$ & \ce{BH3}\\
\rowcolor{primary!5} 3 & \AX{2}{1} & triangulaire & coudée & $<120^\circ$ & \ce{SnCl2}\\
4 & \AX{4}{0} & tétraédrique & tétraédrique & $109{,}5^\circ$ & \ce{CH4}\\
\rowcolor{primary!5} 4 & \AX{3}{1} & tétraédrique & pyramide trigonale & $\approx107^\circ$ & \ce{NH3}\\
4 & \AX{2}{2} & tétraédrique & coudée & $\approx104{,}5^\circ$ & \ce{H2O}\\
\rowcolor{primary!5} 5 & \AX{5}{0} & bipyramide trig. & bipyramide trigonale & $90/120^\circ$ & \ce{PCl5}\\
5 & \AX{4}{1} & bipyramide trig. & à bascule (papillon) & --- & \ce{SF4}\\
\rowcolor{primary!5} 5 & \AX{3}{2} & bipyramide trig. & forme en T & --- & \ce{ClF3}\\
5 & \AX{2}{3} & bipyramide trig. & linéaire & $180^\circ$ & \ce{XeF2}\\
\rowcolor{primary!5} 6 & \AX{6}{0} & octaédrique & octaédrique & $90^\circ$ & \ce{SF6}\\
6 & \AX{5}{1} & octaédrique & pyramide à base carrée & --- & \ce{BrF5}\\
\rowcolor{primary!5} 6 & \AX{4}{2} & octaédrique & plan carré & $90^\circ$ & \ce{XeF4}\\
\bottomrule
\end{tabular}
\end{table}

\begin{astucebox}[le réflexe de fin d'exercice]
Aucun doublet libre sur A ($m=0$) $\Rightarrow$ géométrie moléculaire = géométrie de répulsion.
Dès que $m\geq 1$, on « enlève » les doublets et la forme change (coudée, pyramide, en T,
plan carré\ldots) tandis que l'angle se referme.
\end{astucebox}

\end{document}
